\documentclass[11pt]{article}

% ---------- Packages ----------
\usepackage[margin=1in]{geometry}
\usepackage{amsmath, amssymb, amsthm}
\usepackage{mathtools}
\usepackage{tikz-cd}
% ---------- Theorem environments ----------
\theoremstyle{plain}
\newtheorem{lemma}{Lemma}

% ---------- Common notation ----------
\newcommand{\PP}{\mathbb{P}}
%\newcommand{\mathbb{A}}{\mathbb{A}}
\newcommand{\OO}{\mathcal{O}}
%\newcommand{\Gm}{\mathbb{G}_m}

%\newcommand{\PP}{\mathbb{P}}
\newcommand{\Gm}{\mathbb{G}_{\mathrm{m}}}
\begin{document}

\begin{lemma}[The Euler maps and the Euler exact sequence on $\PP^N$]\label{lem:Euler-sequence-maps}
Let $k$ be a field and let $\PP^N=\PP^N_k$ with homogeneous coordinates
$[x_0:\cdots:x_N]$.  Then there are canonical morphisms of sheaves
\[
\sigma:\OO_{\PP^N}\longrightarrow \OO_{\PP^N}(1)^{\oplus (N+1)},
\qquad
\rho:\OO_{\PP^N}(1)^{\oplus (N+1)}\longrightarrow T_{\PP^N},
\]
with the following properties.

\begin{enumerate}
\item[(i)] \textbf{(Definition of the two maps.)}
The left map $\sigma$ is defined by
\[
\sigma(1)=(x_0,\dots,x_N),
\]
where each $x_i$ is viewed as the standard global section of $\OO_{\PP^N}(1)$.

The right map $\rho$ is defined as follows.  For any open set $U\subset \PP^N$ and
any section $s=(s_0,\dots,s_N)\in \Gamma\!\big(U,\OO_{\PP^N}(1)^{\oplus(N+1)}\big)$,
choose (locally on $U$) homogeneous degree-$1$ representatives of the $s_i$ on the
cone $\mathbb{A}^{N+1}\setminus\{0\}$.  Form the $k$--derivation
\[
\widetilde{X}_s \;:=\; \sum_{i=0}^N s_i\,\frac{\partial}{\partial x_i}.
\]
The induced derivation on $\Gm$--invariant functions (equivalently, on the structure
sheaf of $U$) defines a vector field on $U$; we set $\rho(s)$ to be this vector field.
Equivalently, $\rho(s)$ is the vector field on $\PP^N$ induced by $\widetilde{X}_s$
after quotienting by the $\Gm$--scaling direction.

\item[(ii)] \textbf{(Local coordinate formula for $\rho$.)}
On the standard affine chart $U_0=\{x_0\neq 0\}$, set
\[
y_i=\frac{x_i}{x_0}\qquad (i=1,\dots,N).
\]
Then $\OO_{\PP^N}(1)$ is trivialized on $U_0$ by $x_0$, hence any local section
$s_i\in \Gamma(U_0,\OO_{\PP^N}(1))$ can be written uniquely as $s_i=x_0 f_i$ with
$f_i\in \OO_{\PP^N}(U_0)$.  For $s=(x_0 f_0,\dots,x_0 f_N)$ one has on $U_0$
\[
\rho(s)
\;=\;
\sum_{i=1}^N \bigl(f_i-f_0\,y_i\bigr)\,\frac{\partial}{\partial y_i}.
\]
In particular, for $n=1$ (so $y=x_1/x_0$) one has
\[
\rho(x_0 f_0,\ x_0 f_1) \;=\; (f_1-f_0 y)\,\frac{\partial}{\partial y}.
\]
In particular, when $n=2$ (so $y_1=x_1/x_0$ and $y_2=x_2/x_0$) one has
\[
\rho(x_0 f_0,\ x_0 f_1,\ x_0 f_2)
\;=\;
\bigl(f_1-f_0 y_1\bigr)\,\frac{\partial}{\partial y_1}
\;+\;
\bigl(f_2-f_0 y_2\bigr)\,\frac{\partial}{\partial y_2}.
\]
\item[(iii)] \textbf{(Kernel equals image; Euler exact sequence.)}
On each standard affine chart $U_j=\{x_j\neq 0\}$ (in particular on $U_0$) one has
\[
\ker\!\bigl(\rho|_{U_j}\bigr)=\mathrm{im}\!\bigl(\sigma|_{U_j}\bigr),
\]
and $\sigma$ is injective.  Consequently the sequence
\[
0\longrightarrow \OO_{\PP^N}
\xrightarrow{\ \sigma\ }
\OO_{\PP^N}(1)^{\oplus(N+1)}
\xrightarrow{\ \rho\ }
T_{\PP^N}
\longrightarrow 0
\]
is exact; this is the Euler sequence on $\PP^N$.
\end{enumerate}
\end{lemma}

\begin{proof}
(ii) The definition of $\rho$ is the standard construction: a homogeneous degree-$1$
vector field on the cone descends to a vector field on the projective quotient, and
the Euler vector field $E$ generates the $\Gm$--direction, so it maps to zero on $\PP^n$.
Since $\sigma(1)=(x_0,\dots,x_n)$ corresponds to $E$, we have $\mathrm{im}(\sigma)\subseteq
\ker(\rho)$, and exactness of \eqref{eq:Euler-seq} identifies $\ker(\rho)=\mathrm{im}(\sigma)$.

(2) We compute $\rho(s)$ by evaluating the induced derivation on the affine coordinate
functions $y_i=x_i/x_0$.  Let $\widetilde{X}=\sum_{j=0}^n s_j\,\partial/\partial x_j$
be the corresponding derivation on the coordinate ring of the cone.

\smallskip
\noindent\emph{Step 1: Differentiating an inverse.}
On $U_0$ the function $x_0$ is invertible.  Using the Leibniz rule on
$x_0\cdot x_0^{-1}=1$, we get
\[
0=\widetilde{X}(1)=\widetilde{X}(x_0\cdot x_0^{-1})
=\widetilde{X}(x_0)\,x_0^{-1}+x_0\,\widetilde{X}(x_0^{-1}),
\]
hence
\begin{equation}\label{eq:inverse-rule}
\widetilde{X}(x_0^{-1})=-\,\frac{\widetilde{X}(x_0)}{x_0^2}.
\end{equation}


\noindent\emph{Step 2: Differentiating a quotient (quotient rule).}
For any functions $f,g$ with $g$ invertible, write $f/g=f\cdot g^{-1}$.  Then
\[
\widetilde{X}\!\left(\frac{f}{g}\right)
=\widetilde{X}(f\cdot g^{-1})
=\widetilde{X}(f)\,g^{-1}+f\,\widetilde{X}(g^{-1})
=\frac{\widetilde{X}(f)}{g}-f\,\frac{\widetilde{X}(g)}{g^2},
\]
where we used \eqref{eq:inverse-rule} in the last step.  Thus
\begin{equation}\label{eq:quotient-rule}
\widetilde{X}\!\left(\frac{f}{g}\right)
=\frac{\widetilde{X}(f)\,g-f\,\widetilde{X}(g)}{g^2}.
\end{equation}

\smallskip
\noindent\emph{Step 3: Apply the quotient rule to $y_i=x_i/x_0$.}
Taking $f=x_i$ and $g=x_0$ in \eqref{eq:quotient-rule} gives
\[
\widetilde{X}(y_i)
=\widetilde{X}\!\left(\frac{x_i}{x_0}\right)
=\frac{\widetilde{X}(x_i)\,x_0-x_i\,\widetilde{X}(x_0)}{x_0^2}.
\]
Since $\widetilde{X}=\sum s_j\partial/\partial x_j$, we have
$\widetilde{X}(x_i)=s_i$ and $\widetilde{X}(x_0)=s_0$, hence
\[
\widetilde{X}(y_i)=\frac{s_i x_0-x_i s_0}{x_0^2}.
\]
Now write $s_i=x_0 f_i$ and $x_i=x_0 y_i$.  Substituting yields
\[
\widetilde{X}(y_i)=\frac{(x_0 f_i)x_0-(x_0 y_i)(x_0 f_0)}{x_0^2}
=f_i-f_0 y_i.
\]

\smallskip
\noindent\emph{Step 4: Identify the induced vector field on $U_0$.}
A vector field on $U_0\simeq \mathbb{A}^n$ is determined by its values on the coordinate
functions $y_1,\dots,y_n$, and the derivation sending $y_i\mapsto f_i-f_0 y_i$
is exactly $\sum_{i=1}^n (f_i-f_0 y_i)\,\partial/\partial y_i$. This proves (ii). \\

\begin{proof}
\noindent \textit{Proof of $(3)$.}
Exactness is local on $\PP^N$, so it suffices to work on the standard affine charts
$U_j=\{x_j\neq 0\}$.  We carry out the argument on $U_0$; the other charts are identical
after relabeling coordinates.

\medskip
\noindent\textbf{Step 1: $\sigma$ is injective.}
Since $\OO_{\PP^N}$ is torsion-free and $\sigma$ is nonzero (indeed $\sigma(1)=(x_0,\dots,x_N)$),
it is injective on each affine chart; equivalently, on $U_0$ the component $x_0$ is
invertible, so if $\sigma(f)=0$ then $f x_0=0$ in $\OO(U_0)$, hence $f=0$.

\medskip
\noindent\textbf{Step 2: $\rho\circ\sigma=0$.}
On $U_0$ set $y_i=x_i/x_0$ $(i=1,\dots,N)$.  Write $\sigma(1)=(x_0,\dots,x_N)$ in the
form $(x_0 f_0,\dots,x_0 f_N)$ with $f_0=1$ and $f_i=y_i$ for $i\ge 1$.  Using the local
formula from part~(ii),
\[
\rho(\sigma(1))
=
\sum_{i=1}^N (f_i-f_0 y_i)\,\frac{\partial}{\partial y_i}
=
\sum_{i=1}^N (y_i-y_i)\,\frac{\partial}{\partial y_i}
=0.
\]
Thus $\mathrm{im}(\sigma)\subseteq \ker(\rho)$.

\medskip
\noindent\textbf{Step 3: $\ker(\rho)\subseteq \mathrm{im}(\sigma)$ on $U_0$.}
Let $s=(s_0,\dots,s_N)\in \Gamma\!\big(U_0,\OO_{\PP^N}(1)^{\oplus(N+1)}\big)$ and write
$s_i=x_0 f_i$ with $f_i\in \OO(U_0)$.  If $\rho(s)=0$, then by the local formula
\[
0=\rho(s)=\sum_{i=1}^N (f_i-f_0 y_i)\,\frac{\partial}{\partial y_i}.
\]
Since $\{\partial/\partial y_i\}_{i=1}^N$ are $\OO(U_0)$--linearly independent, we obtain
\[
f_i-f_0 y_i=0 \qquad (i=1,\dots,N),
\]
hence $f_i=f_0 y_i$ for all $i\ge 1$.  Setting $g:=f_0\in \OO(U_0)$ gives
\[
s_0=x_0 f_0=x_0 g=g x_0,
\qquad
s_i=x_0 f_i=x_0(g y_i)=g(x_0 y_i)=g x_i \quad (i\ge 1).
\]
Therefore $s=g(x_0,\dots,x_N)=g\,\sigma(1)$, so $s\in \mathrm{im}(\sigma)$ on $U_0$.
Thus $\ker(\rho|_{U_0})=\mathrm{im}(\sigma|_{U_0})$.

\medskip
\noindent\textbf{Step 4: Surjectivity of $\rho$.}
On $U_0\simeq \mathbb{A}^N$ every vector field has the form
$\sum_{i=1}^N a_i\,\partial/\partial y_i$ with $a_i\in \OO(U_0)$.  Take the section
$s=(0,\ x_0 a_1,\dots,x_0 a_N)$ of $\OO_{\PP^N}(1)^{\oplus(N+1)}$ on $U_0$; then $f_0=0$
and $f_i=a_i$ for $i\ge 1$, so by the local formula
\[
\rho(s)=\sum_{i=1}^N (a_i-0\cdot y_i)\,\frac{\partial}{\partial y_i}
=\sum_{i=1}^N a_i\,\frac{\partial}{\partial y_i}.
\]
Hence $\rho$ is surjective on $U_0$, and therefore surjective globally.

\medskip
Combining Steps~1--4, we conclude that
\[
0\longrightarrow \OO_{\PP^N}
\xrightarrow{\ \sigma\ }
\OO_{\PP^N}(1)^{\oplus(N+1)}
\xrightarrow{\ \rho\ }
T_{\PP^N}
\longrightarrow 0
\]
is exact.
\end{proof}



\end{proof}






\begin{lemma}[The Euler maps and the Euler exact sequence on $\PP^N$]\label{lem:Euler-sequence-maps}
Let $k$ be a field and let $\PP^N=\PP^N_k$ with homogeneous coordinates
$[x_0:\cdots:x_N]$.  Then there are canonical morphisms of sheaves
\[
\sigma:\OO_{\PP^N}\longrightarrow \OO_{\PP^N}(1)^{\oplus (N+1)},
\qquad
\rho:\OO_{\PP^N}(1)^{\oplus (N+1)}\longrightarrow T_{\PP^N},
\]
with the following properties.

\begin{enumerate}
\item[\textup{(i)}] \textbf{(Definition of the two maps.)}
The map $\sigma$ is defined by
\[
\sigma(1)=(x_0,\dots,x_N),
\]
where each $x_i$ is viewed as the standard global section of $\OO_{\PP^N}(1)$.

The map $\rho$ is defined as follows.  For any open set $U\subset \PP^N$ and
any section $s=(s_0,\dots,s_N)\in \Gamma\!\big(U,\OO_{\PP^N}(1)^{\oplus(N+1)}\big)$,
choose (locally on $U$) homogeneous degree-$1$ representatives of the $s_i$ on the
cone $\mathbb{A}^{N+1}\setminus\{0\}$.  Form the $k$-derivation
\[
\widetilde{X}_s \;:=\; \sum_{i=0}^N s_i\,\frac{\partial}{\partial x_i}.
\]
The induced derivation on $\Gm$-invariant functions defines a vector field on $U$;
we set $\rho(s)$ to be this vector field.

Both maps are well defined: $\sigma$ is manifestly a sheaf morphism, and $\rho$ is
independent of the choice of homogeneous representatives because any two
choices differ by a multiple of the Euler vector field
$E=\sum x_i\,\partial/\partial x_i$, which acts as zero on degree-$0$ functions.

\item[\textup{(ii)}] \textbf{(Local coordinate formula for $\rho$.)}
On the standard affine chart $U_0=\{x_0\neq 0\}$, set
\[
y_i=\frac{x_i}{x_0}\qquad (i=1,\dots,N).
\]
Then $\OO_{\PP^N}(1)$ is trivialized on $U_0$ by~$x_0$, so any section
$s_i\in \Gamma(U_0,\OO_{\PP^N}(1))$ can be written uniquely as $s_i=x_0 f_i$ with
$f_i\in \OO_{\PP^N}(U_0)$.  For $s=(x_0 f_0,\dots,x_0 f_N)$ one has
\begin{equation}\label{eq:local-formula}
\rho(s)
\;=\;
\sum_{i=1}^N \bigl(f_i-f_0\,y_i\bigr)\,\frac{\partial}{\partial y_i}.
\end{equation}

\item[\textup{(iii)}] \textbf{(Euler exact sequence.)}
The sequence
\[
0\longrightarrow \OO_{\PP^N}
\xrightarrow{\ \sigma\ }
\OO_{\PP^N}(1)^{\oplus(N+1)}
\xrightarrow{\ \rho\ }
T_{\PP^N}
\longrightarrow 0
\]
is exact.
\end{enumerate}
\end{lemma}

\begin{proof}
The Euler exact sequence is classical; see Hartshorne~\cite[Theorem~II.8.13]{Hartshorne}
for the dual (cotangent) version.  We include a self-contained proof via the
tangent formulation because the local coordinate formula~\eqref{eq:local-formula}
will be used later.

\medskip
\noindent\textbf{Part~(i)} is addressed in the statement; the only non-obvious
point is that $\rho$ is well defined, which follows from the fact that the Euler
vector field $E=\sum x_i\,\partial/\partial x_i$ annihilates every degree-$0$
rational function on $\mathbb{A}^{N+1}\setminus\{0\}$.

\medskip
\noindent\textbf{Part~(ii).}
We compute $\rho(s)$ on $U_0$ by evaluating the derivation
$\widetilde{X}=\sum_{j=0}^N s_j\,\partial/\partial x_j$
on the affine coordinate functions $y_i=x_i/x_0$.
By the quotient rule for derivations,
\[
\widetilde{X}(y_i)
=\widetilde{X}\!\left(\frac{x_i}{x_0}\right)
=\frac{\widetilde{X}(x_i)\,x_0 - x_i\,\widetilde{X}(x_0)}{x_0^2}.
\]
Since $\widetilde{X}(x_i)=s_i$ and $\widetilde{X}(x_0)=s_0$, we get
\[
\widetilde{X}(y_i)=\frac{s_i\, x_0-x_i\, s_0}{x_0^2}.
\]
Substituting $s_i=x_0 f_i$ and $x_i=x_0 y_i$ yields
\[
\widetilde{X}(y_i)
=\frac{(x_0 f_i)\,x_0-(x_0 y_i)(x_0 f_0)}{x_0^2}
=f_i-f_0\, y_i.
\]
A vector field on $U_0\simeq \mathbb{A}^N$ is determined by its values on
$y_1,\dots,y_N$, so
$\rho(s)=\sum_{i=1}^N (f_i-f_0\, y_i)\,\partial/\partial y_i$.

\medskip
\noindent\textbf{Part~(iii).}
Exactness is local, so it suffices to work on each standard chart
$U_j=\{x_j\neq 0\}$.  We treat $U_0$; the other charts are identical
after relabeling.

\smallskip
\emph{Injectivity of $\sigma$.}\;
On $U_0$ the component $x_0$ is invertible, so $\sigma(f)=f(x_0,\dots,x_N)=0$
forces $f\,x_0=0$ in $\OO(U_0)$, hence $f=0$.

\smallskip
\emph{$\rho\circ\sigma=0$.}\;
Write $\sigma(1)=(x_0,\dots,x_N)$ as $(x_0 f_0,\dots,x_0 f_N)$ with $f_0=1$
and $f_i=y_i$ for $i\ge 1$.  By~\eqref{eq:local-formula},
\[
\rho(\sigma(1))
=\sum_{i=1}^N (y_i-1\cdot y_i)\,\frac{\partial}{\partial y_i}=0.
\]

\smallskip
\emph{$\ker(\rho)\subseteq \operatorname{im}(\sigma)$ on $U_0$.}\;
Let $s=(x_0 f_0,\dots,x_0 f_N)$ satisfy $\rho(s)=0$.
By~\eqref{eq:local-formula} and the linear independence of
$\{\partial/\partial y_i\}$, we get $f_i=f_0\,y_i$ for all $i\ge 1$.
Setting $g:=f_0$ gives
\[
s_0=g\,x_0,\qquad s_i=g\,(x_0 y_i)=g\,x_i \quad(i\ge 1),
\]
so $s=g\,(x_0,\dots,x_N)=\sigma(g)$.

\smallskip
\emph{Surjectivity of $\rho$.}\;
Every vector field on $U_0\simeq\mathbb{A}^N$ has the form
$\sum_{i=1}^N a_i\,\partial/\partial y_i$.  Taking $s=(0,\,x_0 a_1,\dots,x_0 a_N)$
gives $f_0=0$ and $f_i=a_i$, so $\rho(s)=\sum a_i\,\partial/\partial y_i$ by
\eqref{eq:local-formula}.
\end{proof}

\begin{thebibliography}{9}
\bibitem{Hartshorne}
R.~Hartshorne,
\emph{Algebraic Geometry},
Graduate Texts in Mathematics~52,
Springer-Verlag, New York, 1977.
\end{thebibliography}

\end{document}


\begin{lemma}[Commutativity of the right square for the $2$--uple Veronese]\label{lem:veronese-right-square}
Let $\nu_2:\PP^1\to \PP^2$ be the second Veronese embedding
\[
\nu_2([u:v])=[U_0:U_1:U_2]=[u^2:uv:v^2].
\]
Let $\rho_{\PP^1}:\OO_{\PP^1}(1)^{\oplus 2}\to T_{\PP^1}$ be the Euler map on $\PP^1$ and
let $\rho_{\PP^2}:\OO_{\PP^2}(1)^{\oplus 3}\to T_{\PP^2}$ be the Euler map on $\PP^2$.
Pulling back the Euler sequence of $\PP^2$ along $\nu_2$ yields a map
\[
\nu_2^*\rho_{\PP^2}:\OO_{\PP^1}(2)^{\oplus 3}\longrightarrow T_{\PP^2}\big|_{\nu_2(\PP^1)}.
\]
Define a morphism of vector bundles
\[
\alpha:\OO_{\PP^1}(1)^{\oplus 2}\longrightarrow \OO_{\PP^1}(2)^{\oplus 3}
\]
by the Jacobian rule
\[
\alpha(a,b)=(2u\,a,\;u\,b+v\,a,\;2v\,b),
\]
and let $\beta=d\nu_2:T_{\PP^1}\to T_{\PP^2}|_{\nu_2(\PP^1)}$ be the differential of $\nu_2$.
Then the following square commutes:
\[
\begin{tikzcd}
\OO_{\PP^1}(1)^{\oplus 2} \arrow[r, "\rho_{\PP^1}"] \arrow[d, "\alpha"'] &
T_{\PP^1} \arrow[d, "\beta=d\nu_2"] \\
\OO_{\PP^1}(2)^{\oplus 3} \arrow[r, "\nu_2^*\rho_{\PP^2}"'] &
T_{\PP^2}\big|_{\nu_2(\PP^1)}
\end{tikzcd}
\]

Equivalently, $\nu_2^*\rho_{\PP^2}\circ \alpha=\beta\circ \rho_{\PP^1}$.
\end{lemma}

\begin{proof}
Since the claim is local on $\PP^1$, it suffices to check it on the affine chart
$U=\{v\neq 0\}\subset \PP^1$.  Set $y=u/v$ on $U$.

\medskip
\noindent\textbf{Step 1: Local coordinates on $\PP^2$ along $\nu_2(U)$.}
On the affine chart $\{U_2\neq 0\}\subset \PP^2$ set
\[
z_0=\frac{U_0}{U_2},\qquad z_1=\frac{U_1}{U_2}.
\]
Along $\nu_2(U)$ we have $U_0=u^2$, $U_1=uv$, $U_2=v^2$, hence
\[
z_0=\frac{u^2}{v^2}=y^2,\qquad z_1=\frac{uv}{v^2}=y.
\]
Therefore the differential $\beta=d\nu_2$ satisfies
\begin{equation}\label{eq:beta-local}
\beta\!\left(\frac{\partial}{\partial y}\right)
=
\frac{d z_0}{d y}\frac{\partial}{\partial z_0}+\frac{d z_1}{d y}\frac{\partial}{\partial z_1}
=
2y\,\frac{\partial}{\partial z_0}+\frac{\partial}{\partial z_1}.
\end{equation}

\medskip
\noindent\textbf{Step 2: Local formula for the Euler maps.}
On $U=\{v\neq 0\}$ the line bundle $\OO_{\PP^1}(1)$ is trivialized by $v$, so any local
section of $\OO_{\PP^1}(1)^{\oplus 2}$ can be written as $(vA,\;vB)$ with
$A,B\in \OO_{\PP^1}(U)$.  By the $n=1$ case of the affine formula from Lemma~\ref{lem:Euler-sequence-maps}, we have
\begin{equation}\label{eq:rhoP1-local}
\rho_{\PP^1}(vA,\;vB)=(A-By)\,\frac{\partial}{\partial y}.
\end{equation}

Similarly, on $\{U_2\neq 0\}$ the bundle $\OO_{\PP^2}(1)$ is trivialized by $U_2$.
After pulling back to $U$, we have $U_2=v^2$, so $\OO_{\PP^1}(2)$ is trivialized by $v^2$.
Hence any local section of $\OO_{\PP^1}(2)^{\oplus 3}$ can be written as
$(v^2C_0,\;v^2C_1,\;v^2C_2)$ with $C_i\in \OO_{\PP^1}(U)$.  The $n=2$ affine formula
gives
\begin{equation}\label{eq:rhoP2-local}
(\nu_2^*\rho_{\PP^2})(v^2C_0,\;v^2C_1,\;v^2C_2)
=
(C_0-C_2 z_0)\frac{\partial}{\partial z_0}+(C_1-C_2 z_1)\frac{\partial}{\partial z_1}.
\end{equation}
Restricting to $\nu_2(U)$, we may substitute $z_0=y^2$ and $z_1=y$.

\medskip
\noindent\textbf{Step 3: The Jacobian map $\alpha$ in local trivializations.}
Let $(vA,\;vB)$ be a local section of $\OO_{\PP^1}(1)^{\oplus 2}$ on $U$.
Using $u=yv$ we compute
\[
\alpha(vA,\;vB)
=
(2u\cdot vA,\; u\cdot vB+v\cdot vA,\;2v\cdot vB)
=
\bigl(v^2(2yA),\;v^2(A+yB),\;v^2(2B)\bigr).
\]
Thus in the notation of \eqref{eq:rhoP2-local},
\begin{equation}\label{eq:C-coeffs}
C_0=2yA,\qquad C_1=A+yB,\qquad C_2=2B.
\end{equation}

\medskip
\noindent\textbf{Step 4: Compare both compositions.}
First compute $(\nu_2^*\rho_{\PP^2})\circ \alpha$ using
\eqref{eq:rhoP2-local} and \eqref{eq:C-coeffs}, and then substitute $z_0=y^2$, $z_1=y$:
\[
(\nu_2^*\rho_{\PP^2})(\alpha(vA,vB))
=
(C_0-C_2 y^2)\frac{\partial}{\partial z_0}+(C_1-C_2 y)\frac{\partial}{\partial z_1}
=
\bigl(2yA-2B y^2\bigr)\frac{\partial}{\partial z_0}+\bigl(A+yB-2By\bigr)\frac{\partial}{\partial z_1}.
\]
Hence
\[
(\nu_2^*\rho_{\PP^2})(\alpha(vA,vB))
=
2y(A-By)\frac{\partial}{\partial z_0}+(A-By)\frac{\partial}{\partial z_1}
=
(A-By)\left(2y\frac{\partial}{\partial z_0}+\frac{\partial}{\partial z_1}\right).
\]
On the other hand, using \eqref{eq:rhoP1-local} and then \eqref{eq:beta-local},
\[
\beta\bigl(\rho_{\PP^1}(vA,vB)\bigr)
=
\beta\!\left((A-By)\frac{\partial}{\partial y}\right)
=
(A-By)\,\beta\!\left(\frac{\partial}{\partial y}\right)
=
(A-By)\left(2y\frac{\partial}{\partial z_0}+\frac{\partial}{\partial z_1}\right).
\]
The two expressions agree, proving $\nu_2^*\rho_{\PP^2}\circ \alpha=\beta\circ \rho_{\PP^1}$
on $U$.  Since the assertion is local and $U$ covers $\PP^1$ together with the chart
$\{u\neq 0\}$, the square commutes globally.
\end{proof}


\begin{lemma}[Commutativity of the right square for the $2$--uple Veronese in $\PP^M$]\label{lem:veronese-right-square-PM}
Let $M\ge 2$ and let $i:\PP^1\to \PP^M$ be a morphism whose image is a smooth conic.
Equivalently, after choosing homogeneous coordinates $[X_0:\dots,X_M]$ on $\PP^M$,
the map $i$ factors as
\[
\PP^1 \xrightarrow{\ \nu_2\ } \PP^2 \xrightarrow{\ \iota\ } \PP^M,
\qquad
\nu_2([u:v])=[u^2:uv:v^2],\quad
\iota([U_0:U_1:U_2])=[U_0:U_1:U_2:0:\cdots:0].
\]
Let $\rho_{\PP^1}:\OO_{\PP^1}(1)^{\oplus 2}\to T_{\PP^1}$ be the Euler map on $\PP^1$ and
let $\rho_{\PP^M}:\OO_{\PP^M}(1)^{\oplus (M+1)}\to T_{\PP^M}$ be the Euler map on $\PP^M$.
Pulling back the Euler sequence of $\PP^M$ along $i$ yields a map
\[
i^*\rho_{\PP^M}:\OO_{\PP^1}(2)^{\oplus (M+1)}\longrightarrow T_{\PP^M}\big|_{i(\PP^1)}.
\]

Define a morphism of vector bundles
\[
\widetilde\alpha:\OO_{\PP^1}(1)^{\oplus 2}\longrightarrow \OO_{\PP^1}(2)^{\oplus (M+1)}
\]
by the Jacobian rule in the first three coordinates and zero in the remaining ones:
\[
\widetilde\alpha(a,b)=\bigl(2u\,a,\;u\,b+v\,a,\;2v\,b,\;0,\dots,0\bigr),
\]
and let $\beta=di:T_{\PP^1}\to T_{\PP^M}|_{i(\PP^1)}$ be the differential of $i$.
Then the following square commutes:
\[
\begin{tikzcd}
\OO_{\PP^1}(1)^{\oplus 2} \arrow[r, "\rho_{\PP^1}"] \arrow[d, "\widetilde\alpha"'] &
T_{\PP^1} \arrow[d, "\beta=di"] \\
\OO_{\PP^1}(2)^{\oplus (M+1)} \arrow[r, "i^*\rho_{\PP^M}"'] &
T_{\PP^M}\big|_{i(\PP^1)}.
\end{tikzcd}
\]
Equivalently, $i^*\rho_{\PP^M}\circ \widetilde\alpha=\beta\circ \rho_{\PP^1}$.
\end{lemma}

\begin{proof}
Since the claim is local on $\PP^1$, it suffices to check it on the affine chart
$U=\{v\neq 0\}\subset \PP^1$.  Set $y=u/v$ on $U$.

\medskip
\noindent\textbf{Step 1: Local coordinates on $\PP^M$ along $i(U)$.}
On the affine chart $\{X_2\neq 0\}\subset \PP^M$ set
\[
z_0=\frac{X_0}{X_2},\qquad z_1=\frac{X_1}{X_2},\qquad w_j=\frac{X_j}{X_2}\ (j=3,\dots,M).
\]
Along $i(U)$ we have $X_0=u^2$, $X_1=uv$, $X_2=v^2$, and $X_j=0$ for $j\ge 3$, hence
\[
z_0=\frac{u^2}{v^2}=y^2,\qquad z_1=\frac{uv}{v^2}=y,\qquad w_j=0\ (j=3,\dots,M).
\]
Therefore the differential $\beta=di$ satisfies
\begin{equation}\label{eq:beta-local-PM}
\beta\!\left(\frac{\partial}{\partial y}\right)
=
\frac{d z_0}{d y}\frac{\partial}{\partial z_0}
+\frac{d z_1}{d y}\frac{\partial}{\partial z_1}
+\sum_{j=3}^M \frac{d w_j}{d y}\frac{\partial}{\partial w_j}
=
2y\,\frac{\partial}{\partial z_0}+\frac{\partial}{\partial z_1},
\end{equation}
since $w_j\equiv 0$ on $i(U)$.

\medskip
\noindent\textbf{Step 2: Local formula for the Euler maps.}
On $U=\{v\neq 0\}$ the line bundle $\OO_{\PP^1}(1)$ is trivialized by $v$, so any local
section of $\OO_{\PP^1}(1)^{\oplus 2}$ can be written as $(vA,\;vB)$ with
$A,B\in \OO_{\PP^1}(U)$.  By the $n=1$ affine formula from Lemma~\ref{lem:Euler-sequence-maps}, we have
\begin{equation}\label{eq:rhoP1-local-PM}
\rho_{\PP^1}(vA,\;vB)=(A-By)\,\frac{\partial}{\partial y}.
\end{equation}

Similarly, on $\{X_2\neq 0\}$ the bundle $\OO_{\PP^M}(1)$ is trivialized by $X_2$.
After pulling back to $U$, we have $X_2=v^2$, so $\OO_{\PP^1}(2)$ is trivialized by $v^2$.
Hence any local section of $\OO_{\PP^1}(2)^{\oplus (M+1)}$ can be written as
\[
(v^2C_0,\;v^2C_1,\;v^2C_2,\;v^2D_3,\dots,v^2D_M)
\]
with $C_i,D_j\in \OO_{\PP^1}(U)$.  The affine Euler formula on the chart $\{X_2\neq 0\}$
gives
\begin{equation}\label{eq:rhoPM-local}
(i^*\rho_{\PP^M})(v^2C_0,\;v^2C_1,\;v^2C_2,\;v^2D_3,\dots,v^2D_M)
=
(C_0-C_2 z_0)\frac{\partial}{\partial z_0}
+(C_1-C_2 z_1)\frac{\partial}{\partial z_1}
+\sum_{j=3}^M (D_j-C_2 w_j)\frac{\partial}{\partial w_j}.
\end{equation}
Restricting to $i(U)$, we may substitute $z_0=y^2$, $z_1=y$, and $w_j=0$.

\medskip
\noindent\textbf{Step 3: The Jacobian map $\widetilde\alpha$ in local trivializations.}
Let $(vA,\;vB)$ be a local section of $\OO_{\PP^1}(1)^{\oplus 2}$ on $U$.
Using $u=yv$ we compute
\[
\widetilde\alpha(vA,\;vB)
=
\bigl(2u\cdot vA,\; u\cdot vB+v\cdot vA,\;2v\cdot vB,\;0,\dots,0\bigr)
=
\bigl(v^2(2yA),\;v^2(A+yB),\;v^2(2B),\;0,\dots,0\bigr).
\]
Thus in the notation of \eqref{eq:rhoPM-local},
\begin{equation}\label{eq:C-D-coeffs}
C_0=2yA,\qquad C_1=A+yB,\qquad C_2=2B,\qquad D_j=0\ (j=3,\dots,M).
\end{equation}

\medskip
\noindent\textbf{Step 4: Compare both compositions.}
First compute $(i^*\rho_{\PP^M})\circ \widetilde\alpha$ using
\eqref{eq:rhoPM-local} and \eqref{eq:C-D-coeffs}, and then substitute
$z_0=y^2$, $z_1=y$, and $w_j=0$:
\begin{align*}
(i^*\rho_{\PP^M})(\widetilde\alpha(vA,vB))
&=
(C_0-C_2 y^2)\frac{\partial}{\partial z_0}
+(C_1-C_2 y)\frac{\partial}{\partial z_1}
+\sum_{j=3}^M (D_j-C_2\cdot 0)\frac{\partial}{\partial w_j}\\
&=
\bigl(2yA-2B y^2\bigr)\frac{\partial}{\partial z_0}
+\bigl(A+yB-2By\bigr)\frac{\partial}{\partial z_1}
+\sum_{j=3}^M 0\cdot \frac{\partial}{\partial w_j}\\
&=
(A-By)\left(2y\frac{\partial}{\partial z_0}+\frac{\partial}{\partial z_1}\right).
\end{align*}
On the other hand, using \eqref{eq:rhoP1-local-PM} and then \eqref{eq:beta-local-PM},
\[
\beta\bigl(\rho_{\PP^1}(vA,vB)\bigr)
=
\beta\!\left((A-By)\frac{\partial}{\partial y}\right)
=
(A-By)\,\beta\!\left(\frac{\partial}{\partial y}\right)
=
(A-By)\left(2y\frac{\partial}{\partial z_0}+\frac{\partial}{\partial z_1}\right).
\]
The two expressions agree, proving $i^*\rho_{\PP^M}\circ \widetilde\alpha=\beta\circ \rho_{\PP^1}$
on $U$.  Since the assertion is local and $U$ covers $\PP^1$ together with the chart
$\{u\neq 0\}$, the square commutes globally.
\end{proof}



\end{document}
